\section{引言}\label{sec:intro}

部分子分布函数（PDF）为以近光速运动的核子内部的夸克与胶子提供了有效的描述~\cite{Ellis:1991qj,Thomas:2001kw}。它们在计算涉及核子的硬散射截面时起着关键作用~\cite{Gao:2017yyd}，其唯象抽取的不确定度一直是大型强子对撞机（LHC）及其他强子装置上理论预言误差的主要来源之一。因此，精确掌握 PDF 无论对标准模型（SM）的精确检验，还是对寻找超出标准模型的新物理都十分重要。另一方面，核子内部分子的密度直接提供核子内在性质的信息，例如核子自旋与质量的起源，以及海夸克在各种物理量中所起的作用~\cite{Accardi:2012qut}。

然而，从量子色动力学（QCD）的第一性原理直接计算部分子物理一直是一项困难的任务。对这方面各种努力的综述可参见文献~\cite{Cichy:2018mum}。大动量有效场理论（LaMET）的提出是应对这一挑战的重要一步~\cite{Ji:2013dva,Ji:2014hxa,Ji:2020ect}。迄今为止，LaMET 已被广泛用于计算夸克同位旋矢量分布函数~\cite{Lin:2014zya,Alexandrou:2015rja,Chen:2016utp,Alexandrou:2016jqi,Alexandrou:2018pbm,Chen:2018xof,Lin:2018pvv,Liu:2018uuj,Alexandrou:2018eet,Liu:2018hxv,Chen:2018fwa,Izubuchi:2019lyk,Shugert:2020tgq,Chai:2020nxw,Lin:2020ssv,Fan:2020nzz}、广义部分子分布~\cite{Chen:2019lcm,Alexandrou:2019dax}、分布振幅（DA）~\cite{Zhang:2017bzy,Chen:2017gck,Zhang:2020gaj}以及横向动量依赖分布~\cite{Shanahan:2019zcq,Shanahan:2020zxr,Zhang:2020dbb}。

LaMET 建议在有限动量的核子中计算与时间无关的物理分布，它们是欧氏观测量。这类有限动量量随后可用有效场论技术匹配到光前（LF）部分子性质~\cite{Ji:2020byp}。例如，对共线夸克分布，建议先在格点上计算欧氏空间关联函数（或称准光前（quasi-LF）关联函数）
\bea\label{eq:quasi}
\tilde{h}(z, P_z) = \langle P| O_{\gamma_t}(z)|P \rangle,
\eea
其中 $|P\rangle$ 是带有大动量 $P$ 的核子态，非定域算符为
\bea\label{eq:operator}
O_{\Gamma}(z)=\bar{\psi}(0) \Gamma U(0,z) \psi (z),
\eea
式中 $\psi$ 与 $\bar{\psi}$ 表示裸夸克场，$\Gamma$ 是一个 Dirac 结构，而 $U(0, z) = \exp(-ig\int_0^{z} dz' A_z(z'))$ 是沿 $z$ 方向的 Wilson 链接，$A^\mu$ 是胶子规范势。在对 $\tilde{h}$ 作恰当的重正化并把它匹配到某个连续方案——如维数正规化与（修正）最小减除 ($\overline{\rm MS}$)——之后，通过傅里叶变换便得到所谓的夸克准 PDF，再经由微扰 QCD 匹配即可提取夸克 PDF。

在格点正规化下重正化准光前关联 $\tilde{h}$ 是一件不平凡的工作。为说明这一点，让我们举一个简单的例子：$O_\Gamma(z)$ 在带有大欧氏动量 $p^2$ 的离壳夸克态 $|q\rangle$ 中的矩阵元（实际上是一个离壳截腿格林函数）。它在单圈水平具有如下形式~\cite{Alexandrou:2017huk}，
\bea\label{eq:1-loop}
\langle O_{\Gamma}(z)\rangle =\Gamma\left(1+\gamma g^2 \textrm{log}(z^2/a^2)-m_{-1}\frac{z}{a}+\cdots\cdots\right),\nonumber\\
\eea
其中 $g$ 是裸规范耦合，$a$ 是格点间距。与对数项系数 $\gamma$ 不同，线性发散系数 $m_{-1}$ 可能对格点上费米子与规范作用量的细节敏感，因为这些作用量本身定义了格点正规化。线性发散项正比于以格点单位计的 Wilson 链接长度 $z/a$；当计入高阶修正后，它在大 $z$ 或小 $a$ 时可以指数级增大。因此，在把格点计算的准 PDF 外推到连续极限 $a\rightarrow0$ 并最终匹配到连续方案的 PDF 之前，必须先消除线性发散效应。

迄今的理论研究表明，在连续理论中准 PDF 算符是可乘性重正化的~\cite{Ji:2015jwa, Ji:2017oey,Ishikawa:2017faj,Green:2017xeu}。在格点上，由于 $z\to 0$ 与 $a\to 0$ 两个极限不对易，最近有工作建议采用混合方案分别重正化大距离与短距离关联~\cite{Ji:2020brr}，其优点是避免某些离散化效应以及重正化阶段引入的非期望的红外效应。在短距离展开适用的小 $z$ 处，可以使用各种比值方案，例如除以正规化无关动量减除（RI/MOM）方案下准 PDF 算符的离壳夸克矩阵元~\cite{Green:2017xeu,Chen:2017mzz,Alexandrou:2017huk}，或除以静止系强子矩阵元 $\tilde{h}(z, P_z=0)$~\cite{Orginos:2017kos,Izubuchi:2018srq}。在大距离处，可以通过减除幂次发散直接去除 Wilson 线自能量的效应~\cite{Chen:2016fxx,Ji:2017oey,Ishikawa:2017faj,Green:2017xeu}。除使用 RI/MOM 因子或静止系强子矩阵元之外，文献还提出了其他提取线性发散的方法，包括 Wilson 圈~\cite{Chen:2016fxx,Zhang:2017bzy}、真空期望值 $\langle O_{\Gamma}(z)\rangle$~\cite{Braun:2018brg} 以及规范固定的 Wilson 链接~\cite{Ji:2020ect} 等~\cite{Ji:2020brr}。

然而，在数值上减除线性发散是一项极其精细的工作。首先，线性发散可能对格点计算中的系统误差敏感：数据里，具有相同线性发散的两个矩阵元之间可能存在细微差异，这些差异可能导致上述方法失效，对小格点间距尤其如此。其次，某些重正化因子还存在别的问题。例如 $\langle O_{\Gamma}(z)\rangle$ 的真空期望值在短距离处的领头贡献为零，因此数值上提取相应的线性发散非常困难。最后，正如我们将看到的，Wilson 费米子线性发散的手征对称性破缺效应可能使线性发散失去普适性~\cite{Zhang:2020rsx}。

为了更好地理解 LaMET 应用中线性和发散带来的影响，我们对用不同格点作用量生成的多组格点数据系统地研究其矩阵元的线性发散。我们提出一种自重正化方法，当拥有若干不同格点间距的数据时，它可以消除全部发散和离散化误差。该方法的思路是直接从待重正化的矩阵元自身提取重正化因子和剩余贡献，而不借助另一个额外的矩阵元来作重正化。随后我们把这样经验性重正化的矩阵元与连续统中的对应量进行匹配。我们的方法在很大程度上不依赖模型，拟合函数的每一项都有物理上的动机。检验表明，该方法对我们考虑的所有数据组都行之有效；这些数据包括用价 clover 与 overlap 作用量、在 MILC~\cite{Bazavov:2012xda} $N_f=2+1+1$ 与 RBC~\cite{Blum:2014tka} $N_f=2+1$ 组态上、格点间距从 $0.03$ fm 到 $0.12$ fm 计算的矩阵元。我们还发现该方法同样适用于经过涂抹处理的矩阵元，而涂抹正是改善统计精度所需要的。

本文其余部分安排如下。第~\ref{sec:theory} 节给出微扰重正化理论。我们所计算矩阵元的定义与模拟设置见第~\ref{sec:setup} 节。第~\ref{sec:test} 节分析不同系综下的线性发散。第~\ref{sec:StrtoRe} 节给出矩阵元自重正化的策略，并汇集在不同态、不同价费米子作用下计算的 $O_{\gamma_t}(z)$ 矩阵元的结果。第~\ref{sec:highorder} 节讨论其他拟合选项，其主要差别在于高阶项的处理。在第~\ref{sec:test}、\ref{sec:StrtoRe} 和~\ref{sec:highorder} 节中，我们只分析未经 HYP 涂抹的矩阵元。第~\ref{sec:smearing} 节把方法推广到 HYP 涂抹情形。第~\ref{sec:otherresult} 节在若干真空态矩阵元（包括 Wilson 圈、准 PDF 算符和规范固定 Wilson 链接）中检验 HYP 涂抹情形的线性发散。第~\ref{sec:summary} 节总结全文并讨论有待进一步研究的问题。
